\scanpage{23}
$S$ \inserted{\unsure{nœth|noeth}} muni d'un Idéal \inserted{\deleted{\unsure{qu}}coh.} de carré nul $J$,

$Z\xrightarrow{\unsure{\varphi,v|v}}X;Y$ sur $S$ d'où $v_0:Z_0\mathbin{-}X_0,Y_0$ sur $S_0$

\[
S_0=V(J)
\]

$u_0:X_0\to Y_0$ tel que $u_0\varphi_0=v_0$

% slanted insertion in left margin, with a leader to the equation below
On va étudier l'obstruction à relever $u_0$ en $u:X\to Y$ tel que $u\varphi=v$\inserted{ et l'obstruction analogue pour la situation déduite à l'aide d'un hom \unsure{$u$|$w$}$:J\to G$, ou après localisation\ldots} À priori, \deleted{l'obstruction à relever}\inserted{les relèvements de} $u_0$ en un $u$ (sans condition supplémentaire imposée) correspondent aux splittings d'une extension sur $P_0$

\[
0\mathbin{-}\unsure{\mathcal O_{\Gamma_0}|\underline{\mathcal O}_{\Gamma_0}}\otimes_{S_0}J\longrightarrow E\longrightarrow\underline I_{\Gamma_0}\mathbin{-}0
\]

% slanted

\margin{SGA 3 \uline{III} 4.2}

les relèvts pour $v_0$ correspondent aux splittings d'une extension \deleted{\unsure{de}} sur $P'_0=Z_0\times_{\unsure{S_0|S}}Y_0$

\[
0\mathbin{-}\mathcal O_{\Gamma'_0}\otimes_{S_0}J\mathbin{-}E'\mathbin{-}\underline I_{\Gamma'_0}\mathbin{-}0
\]

D'ailleurs, on constate qu'on a un homomorphisme d'extensions, compatible avec $P'_0\to P_0$

\diagram{\begin{tikzcd}[ampersand replacement=\&,column sep=small]
0\arrow[r] \& \mathcal O_{\Gamma_0}\otimes_{S_0}J\arrow[r]\arrow[d] \& E\arrow[r]\arrow[d] \& I_{\Gamma_0}\arrow[r]\arrow[d] \& 0\\
0\arrow[r,dash] \& \mathcal O_{\Gamma'_0}\otimes_{S_0}J\arrow[r,dash] \& E'\arrow[r] \& \underline I_{\Gamma'_0}\arrow[r] \& 0
\end{tikzcd}}

(en fait, $E'$ se déduit de $E$ par extension de la base $P'_{\unsure{\xi|0}}\mathbin{-}P_{\unsure{\xi|0}}$ \unsure{ou} avec

\diagram{\begin{tikzcd}[ampersand replacement=\&,column sep=small,row sep=small]
\& P_0\arrow[d]\\ Z_0\arrow[r] \& X_0
\end{tikzcd}}

\unsure{;]|;} il faut préciser ce pt \ldots{} ]. \deleted{\unsure{Considérons, pour un|Conditions, puis un}}\deleted{\unsure{$G$ quasi-cohérent sur}}\inserted{D'ailleurs, la} deuxième ligne est munie d'un splitting, et les solutions cherchées sont les splittings de la première \uline{compatibles} avec celui donné pour la seconde.

% archive sheet number: 22
